\section*{Chapitre \ref{ch:b1:atmospheric-circulation} --- La circulation générale de l'atmosphère}
% values: calc/ch28.py

\begin{solution}{exo:b1:atmospheric-circulation:1}
(a)~$f = 2\Omega\sin\varphi$~: $0$ à l'équateur, \qty{7.29e-5}{s^{-1}} à \ang{30},
\qty{1.031e-4}{s^{-1}} à \ang{45} et \qty{1.458e-4}{s^{-1}} au pôle.
(b)~$fV = \num{1.031e-4} \times 10 = \qty{1.03e-3}{m.s^{-2}}$, environ \num{1.05e-4}
de $g$. (c)~Dans l'hémisphère Nord la déviation se fait à droite du mouvement~:
un vent soufflant vers le nord est poussé vers l'est.
\end{solution}

\begin{solution}{exo:b1:atmospheric-circulation:2}
(a)~$|\nabla p|/\rho = (400/\num{4e5})/1.225 = \qty{8.16e-4}{m.s^{-2}}$.
(b)~$10/\num{8.16e-4} = \qty{12250}{s}$, \qty{3.4}{h}. (c)~L'accélération de Coriolis
d'un vent de \qty{10}{m/s} à \ang{45}~N est de \qty{1.03e-3}{m.s^{-2}}~: les deux sont de
même taille (rapport $0.79$). Quelques heures après le début du mouvement, l'air subit une
force de Coriolis capable d'équilibrer le gradient de pression~: le vent se fixe près de la
vitesse géostrophique, ici $\num{8.16e-4}/\num{1.031e-4} = \qty{7.9}{m/s}$, au lieu
d'accélérer en descendant le gradient.
\end{solution}

\begin{solution}{exo:b1:atmospheric-circulation:3}
(a)~$f = 2 \times \num{7.292e-5} \times \sin\ang{40} = \qty{9.374e-5}{s^{-1}}$~;
$|\vec v_g| = (400/\num{2.5e5})/(1.225 \times \num{9.374e-5}) = \qty{13.9}{m/s}$, soit
$13.93/0.5144 = \qty{27.1}{kn}$. (b)~Dans l'hémisphère Nord le vent garde la
basse à sa gauche~: avec la basse au sud, il souffle vers l'ouest, venant de l'est
(\omterm{def:b1:spherical-astronomy:horizontal}{azimut} \ang{90}). (c)~Le fût pointe vers l'est depuis le cercle de la station, vers le
vent~; \qty{27.1}{kn} s'arrondit à \qty{25}{kn}, deux barbules entières et une demi-barbule,
dessinées du côté horaire du fût, ici son côté sud, vers la basse.
\end{solution}

\begin{solution}{exo:b1:atmospheric-circulation:4}
(a)~\ang{15}~N~: \omterm{def:b1:atmospheric-circulation:winds}{alizés} de nord-est~; \ang{15}~S~: \omterm{def:b1:atmospheric-circulation:winds}{alizés} de sud-est~;
\ang{45}~N et \ang{45}~S~: \omterm{def:b1:atmospheric-circulation:winds}{vents d'ouest}. (b)~À \ang{15} l'air de surface de la \omterm{def:b1:atmospheric-circulation:cells}{cellule de Hadley}
retourne vers l'équateur et est dévié vers l'ouest par la force de Coriolis
(à droite au nord, à gauche au sud)~: vent d'est, $u$ négatif. À \ang{45} l'air de
surface de la \omterm{def:b1:atmospheric-circulation:cells}{cellule de Ferrel} va vers le pôle et est dévié vers l'est~: vent d'ouest,
$u$ positif. (c)~Des trois \omterm{def:b1:spherical-astronomy:geographic}{latitudes}, la pression est la plus élevée à \ang{30}~N,
\qty{1016.8}{hPa}~: l'\omterm{def:b1:atmospheric-circulation:cells}{anticyclone subtropical} se trouve près de \ang{30}~N (le profil complet place
son maximum près de \ang{35}~N).
\end{solution}

\begin{solution}{exo:b1:atmospheric-circulation:5}
(a)~$|\nabla p|/\rho = (200/\num{5e5})/1.225 = \qty{3.27e-4}{m.s^{-2}}$, divisé par
$f$~: \qty{2.6}{m/s} à \ang{60}, \qty{2.9}{m/s} à \ang{50}, \qty{4.5}{m/s} à
\ang{30}, \qty{6.5}{m/s} à \ang{20}, \qty{12.9}{m/s} à \ang{10} et
\qty{25.7}{m/s} à \ang{5}. (b)~La force de Coriolis tourne un vent en un temps de l'ordre de
$1/f$~: \qty{2.2}{h} à \ang{60}, \qty{11}{h} à \ang{10}, \qty{22}{h} à \ang{5}, et
infini à l'équateur. Près de l'équateur le champ de pression change, et le
frottement agit, bien avant que le vent puisse être tourné le long des isobares~: l'équilibre
ne s'installe jamais, et l'air s'écoule en partie en descendant le gradient, à vitesses modestes.
\end{solution}

\begin{solution}{exo:b1:atmospheric-circulation:6}
(a)~$u = 464.6\,\sin^2\varphi/\cos\varphi$~: \qty{14.2}{m/s} à \ang{10},
\qty{32.2}{m/s} à \ang{15}, \qty{57.8}{m/s} à \ang{20}, \qty{134.1}{m/s} à
\ang{30}. (b)~$u = 464.6\,(\cos^2\ang{10} - \cos^2\varphi)/\cos\varphi$~:
\qty{17.7}{m/s} à \ang{15}, \qty{42.9}{m/s} à \ang{20}, \qty{117.9}{m/s} à
\ang{30}. (c)~L'air qui part plus loin de l'équateur part plus près de l'axe,
avec moins de moment cinétique en réserve~: il arrive avec un vent d'ouest plus faible à chaque
\omterm{def:b1:spherical-astronomy:geographic}{latitude}. L'étalement de la branche ascendante abaisse le \omterm{def:b1:atmospheric-circulation:jet}{courant-jet} prévu, mais de loin pas
assez pour atteindre les \qty{40}{m/s} observés.
\end{solution}

\begin{solution}{exo:b1:atmospheric-circulation:7}
(a)~Le vent souffle vers le nord-est. Le dos tourné à lui, face au nord-est,
la basse est à gauche~: au nord-ouest. (b)~Le vent souffle vers le
sud-est~; dans l'hémisphère Sud la basse est à droite d'une personne dont le
dos est tourné au vent~: au sud-ouest. (c)~Le vent de surface est tourné par rapport aux
isobares vers la basse, c'est-à-dire dans le sens trigonométrique par rapport au \omterm{def:b1:atmospheric-circulation:geostrophic}{vent géostrophique} dans
l'hémisphère Nord. Le \omterm{def:b1:atmospheric-circulation:geostrophic}{vent géostrophique}, et la basse à angle droit à sa gauche,
sont donc tournés dans le sens horaire par rapport à la réponse de (a)~: la basse se trouve entre
le nord-ouest et le nord, à l'\omterm{def:b1:spherical-astronomy:horizontal}{azimut} $\ang{315} + \alpha$ pour un angle de coupe
$\alpha$.
\end{solution}

\begin{solution}{exo:b1:atmospheric-circulation:8}
(a)~L'aire entre l'équateur et \ang{30}~N est $2\pi a_\oplus^2\sin\ang{30} =
\qty{1.275e14}{m^2}$~; la masse au-dessus d'un mètre carré entre les deux niveaux est
$\Delta p/g = \num{9e4}/9.81 = \qty{9174}{kg}$~; le total est
\qty{1.170e18}{kg}. (b)~$\num{1.170e18}/\num{8.7e10} = \qty{1.34e7}{s}$, environ
\num{156} jours, cinq mois. (c)~$\num{1.170e18}/\num{21.4e10} = \qty{5.47e6}{s}$,
environ \num{63} jours, deux mois.
\end{solution}

\begin{solution}{exo:b1:atmospheric-circulation:9}
(a)~$\num{0.01} \times \num{86400} = \qty{864}{m}$ par jour. (b)~$0.864 \times 9.8 =
\qty{8.5}{K}$ par jour. (c)~L'air se refroidit en émettant du rayonnement infrarouge~:
l'atmosphère perd plus d'énergie par rayonnement qu'elle n'en absorbe, et la descente de
l'air, qui le réchauffe par compression, compense la perte. (d)~L'air qui descend se réchauffe et s'assèche
par rapport à la saturation, ses nuages s'évaporent et la pluie ne peut se former~: les \omterm{def:b1:atmospheric-circulation:cells}{anticyclones
subtropicaux} sont des ceintures de ciel clair et de peu de pluie, où se forment des déserts chauds sur les terres.
\end{solution}

\begin{solution}{exo:b1:atmospheric-circulation:10}
(a)~$V^2 + frV - fr|\vec v_g| = 0$ avec $fr = \qty{55.85}{m/s}$~:
$V = \tfrac12\bigl(-55.85 + \sqrt{55.85^2 + 4 \times 55.85 \times 15}\bigr) =
\qty{12.3}{m/s}$, \qty{82}{\percent} du \omterm{def:b1:atmospheric-circulation:geostrophic}{vent géostrophique}. (b)~$V^2 - frV +
fr|\vec v_g| = 0$ n'a de racines réelles que si $(fr)^2 \geq 4fr|\vec v_g|$, c'est-à-dire
$|\vec v_g| \leq fr/4 = \qty{14.0}{m/s}$. (c)~Près du centre d'une haute, où $r$ est
petit, la limite $fr/4$ est petite~: le gradient de pression doit être faible, les isobares éloignées
et les vents légers. Une dépression n'a pas de telle limite~: ses isobares peuvent se serrer
près du centre, avec des vents forts.
\end{solution}

\begin{solution}{exo:b1:atmospheric-circulation:11}
(a)~Pour de petits $\varphi$, $\sin\varphi \approx \varphi$ et $\cos\varphi \approx 1$,
donc $u \approx \Omega a_\oplus\varphi^2$. (b)~$\ang{10} = \qty{0.1745}{rad}$~:
$464.6 \times 0.1745^2 = \qty{14.15}{m/s}$, contre \qty{14.23}{m/s} exactement,
soit \qty{0.5}{\percent} de moins. (c)~La résolution de $464.6\sin^2\varphi/\cos\varphi = 27.0$
numériquement donne $\varphi = \ang{13.7}$. (d)~L'anneau égale déjà le \omterm{def:b1:atmospheric-circulation:jet}{courant-jet}
annuel moyen réel à moins de \ang{14}~: la conservation du moment cinétique ne décrit la
branche supérieure de la \omterm{def:b1:atmospheric-circulation:cells}{cellule de Hadley} que dans les tropiques profonds~; au-delà d'environ
\ang{15}, les tourbillons prennent le relais.
\end{solution}

\begin{solution}{exo:b1:atmospheric-circulation:12}
(a)~Moyenne $31.7/12 = \ang{2.6}$~N~; étendue de \ang{12.8}~S à \ang{10.9}~N, soit
\ang{23.7}. (b)~La position la plus au nord arrive en septembre, trois mois après le
solstice de juin~; la plus au sud en février, deux mois après le solstice de décembre.
Le retard vient du stockage de chaleur, surtout dans la \omterm{def:b1:fluid-envelopes-structure:ocean-layers}{couche de mélange} de l'océan
(\cref{ch:b1:fluid-envelopes-structure}), qui se réchauffe et se refroidit lentement. (c)~La moyenne
zonale moyenne toutes les \omterm{def:b1:spherical-astronomy:geographic}{longitudes}~: si sur les continents la bande de pluie oscille loin
dans l'hémisphère d'été, la moyenne de toutes les \omterm{def:b1:spherical-astronomy:geographic}{longitudes} peut traverser l'équateur
bien que la bande sur un océan ne le fasse pas.
\end{solution}

\begin{solution}{pb:b1:atmospheric-circulation:1}
\begin{enumerate}
  \item $\rho = p/(R_dT) = \num{100600}/(287.05 \times 283.15) =
        \qty{1.238}{kg.m^{-3}}$.
  \item $|\nabla p| = \qty{400}{Pa}/\qty{2.5e5}{m} = \qty{1.6e-3}{Pa.m^{-1}}$.
  \item $|\nabla p|/\rho = \num{1.6e-3}/1.238 = \qty{1.293e-3}{m.s^{-2}}$.
  \item $\Omega = 2\pi/\num{86164.2} = \qty{7.292e-5}{rad.s^{-1}}$~;
        $f = 2\Omega\sin\ang{50} = \qty{1.117e-4}{s^{-1}}$.
  \item $|\vec v_g| = \num{1.293e-3}/\num{1.117e-4} = \qty{11.57}{m/s}$.
  \item $11.57/0.514444 = \qty{22.5}{kn}$~; $11.57 \times 3.6 = \qty{41.7}{km/h}$.
  \item Le temps est $|\vec v_g|/(|\nabla p|/\rho)$ et $|\vec v_g| = (|\nabla
        p|/\rho)/f$, donc il vaut $1/f = \qty{8951}{s}$, \qty{2.49}{h}~; et $2\pi/f =
        \qty{15.6}{h}$. Les configurations qui durent un jour ou plus sont donc proches de
        l'équilibre.
  \item La \omterm{def:b1:atmospheric-circulation:pgf}{force de gradient de pression} pointe vers le nord, vers la basse. Le \omterm{def:b1:atmospheric-circulation:geostrophic}{vent
        géostrophique} souffle le long des isobares avec la basse à sa gauche~: vers l'est,
        venant de l'ouest, \omterm{def:b1:spherical-astronomy:horizontal}{azimut} \ang{270}~: un vent d'ouest.
  \item Le dos tourné à un vent d'ouest, face à l'est, la gauche est le nord, où se trouve la basse~:
        la loi est satisfaite.
  \item Le fût pointe vers l'ouest depuis le cercle de la station, vers l'endroit d'où vient le vent~;
        \qty{22.5}{kn} s'arrondit à \qty{20}{kn}, deux barbules entières~; elles se placent du
        côté horaire du fût, son côté nord, vers la basse.
  \item À \ang{50}~S, $f < 0$~: le vent garde la basse à sa droite, donc il souffle
        vers l'ouest, venant de l'est.
  \item $f(\ang{30}) = \qty{7.292e-5}{s^{-1}}$~: \qty{17.7}{m/s}~; $f(\ang{60}) =
        \qty{1.263e-4}{s^{-1}}$~: \qty{10.2}{m/s}.
  \item $f(\ang{5}) = \qty{1.271e-5}{s^{-1}}$ donne \qty{101.7}{m/s}, un vent jamais
        observé avec de tels gradients. Le temps de rotation $1/f$ est \qty{22}{h}~:
        l'équilibre n'a pas le temps de s'installer, et la formule ne s'applique pas.
  \item $9.0/11.57 = 0.778$~; le long des isobares $9.0\cos\ang{25} = \qty{8.16}{m/s}$,
        à travers elles vers la basse $9.0\sin\ang{25} = \qty{3.80}{m/s}$.
  \item Le \omterm{def:b1:atmospheric-circulation:geostrophic}{vent géostrophique} souffle vers l'\omterm{def:b1:spherical-astronomy:horizontal}{azimut} \ang{90}~; tourné de \ang{25} vers
        la basse au nord, le vent de surface souffle vers \ang{65}, venant de
        \ang{245}, l'ouest-sud-ouest.
  \item $\tan\alpha = 1/(f\tau)$ donne $\tau = 1/(\num{1.117e-4}\tan\ang{25}) =
        \qty{1.92e4}{s}$, \qty{5.3}{h}~; le modèle prédit $11.57\cos\ang{25} =
        \qty{10.5}{m/s}$, plus que les \qty{9.0}{m/s} observés~: le frottement réel
        n'est pas exactement linéaire en fonction du vent, et le ralentit davantage que le modèle à cet
        angle.
  \item À travers un cylindre de rayon $R = \qty{500}{km}$ et de hauteur $h =
        \qty{1}{km}$~: $\dot m = 2\pi R h\rho u_r = 2\pi \times \num{5e5} \times 1000
        \times 1.238 \times 3.80 = \qty{1.48e10}{kg/s}$.
  \item $w = \dot m/(\rho\pi R^2) = 2hu_r/R = 2 \times 1000 \times 3.80/\num{5e5} =
        \qty{0.0152}{m/s}$, environ \qty{1310}{m} par jour.
  \item $1.31 \times 9.8 = \qty{12.9}{K}$ par jour~: l'air atteint vite la saturation,
        et la dépression se remplit de nuages et de pluie. Sous un \omterm{def:b1:atmospheric-circulation:cyclone}{anticyclone} l'air de
        surface s'écoule vers l'extérieur, l'air au-dessus descend et se réchauffe, et le ciel
        se dégage.
  \item $\Omega a_\oplus = \num{7.292e-5} \times \num{6.371e6} = \qty{464.6}{m/s}$
        à l'équateur~; $464.6\cos\ang{30} = \qty{402.3}{m/s}$ à \ang{30}~N.
  \item $u = 464.6\,\sin^2\varphi/\cos\varphi$~: \qty{57.8}{m/s} à \ang{20}~N et
        \qty{134.1}{m/s} à \ang{30}~N.
  \item $u = 464.6\,(\cos^2\ang{10} - \cos^2\ang{30})/\cos\ang{30} =
        \qty{117.9}{m/s}$.
  \item $134.1/42.3 = 3.17$. Les tourbillons des moyennes \omterm{def:b1:spherical-astronomy:geographic}{latitudes}, dépressions et
        \omterm{def:b1:atmospheric-circulation:cyclone}{anticyclones}, prennent du moment cinétique à la branche supérieure et l'emportent
        ailleurs~; et l'air de la cellule monte sur une bande de \omterm{def:b1:spherical-astronomy:geographic}{latitudes} plutôt
        qu'au repos à l'équateur. La moyenne zonale moyenne aussi un \omterm{def:b1:atmospheric-circulation:jet}{courant-jet} sinueux
        sur toutes les \omterm{def:b1:spherical-astronomy:geographic}{longitudes}.
  \item $M/M_0 = (402.3 + 42.3)\cos\ang{30}/464.6 = 0.829$~: l'air a perdu
        \qty{17}{\percent} du moment cinétique qu'il aurait gardé comme anneau libre.
  \item Le \omterm{def:b1:atmospheric-circulation:geostrophic}{vent géostrophique} en P est de \qty{11.6}{m/s}, venant de l'ouest~; le vent
        mesuré au sol est de \qty{9.0}{m/s}, environ les quatre cinquièmes de lui, et tourné de
        \ang{25} vers la basse par le frottement.
\end{enumerate}
\end{solution}
